% 图 3-1 KSE SoC 总体框图
\documentclass[border=4pt]{standalone}
% 所有框图共用的样式（standalone + TikZ，xelatex/tectonic 编译）
\usepackage{fontspec}
\usepackage{xeCJK}
\setCJKmainfont{Noto Sans CJK SC}[BoldFont=Noto Sans CJK SC Bold]
\setCJKsansfont{Noto Sans CJK SC}[BoldFont=Noto Sans CJK SC Bold]
\renewcommand{\familydefault}{\sfdefault}
\usepackage{tikz}
\usetikzlibrary{positioning,arrows.meta,fit,calc,backgrounds,shapes.geometric,decorations.pathreplacing,patterns}
% 统一配色：灰蓝主色，橙色强调，紫色辅助
\definecolor{cBlue}{HTML}{2E6DA4}
\definecolor{cBlueL}{HTML}{DCE7F2}
\definecolor{cBlueM}{HTML}{A9C3DD}
\definecolor{cOrange}{HTML}{C9732E}
\definecolor{cOrangeL}{HTML}{F6E3D2}
\definecolor{cPurple}{HTML}{8A5BB0}
\definecolor{cPurpleL}{HTML}{ECE3F4}
\definecolor{cGray}{HTML}{5B6670}
\definecolor{cGrayL}{HTML}{EEF0F2}
\definecolor{cInk}{HTML}{25303B}
\tikzset{
  >={Stealth[length=2.2mm,width=1.6mm]},
  every node/.style={font=\small, text=cInk},
  blk/.style={draw=cBlue, fill=cBlueL, line width=0.6pt, rounded corners=1.5pt, align=center, minimum height=7mm, inner sep=3pt},
  blkO/.style={blk, draw=cOrange, fill=cOrangeL},
  blkP/.style={blk, draw=cPurple, fill=cPurpleL},
  blkG/.style={blk, draw=cGray, fill=cGrayL},
  blkW/.style={blk, draw=cGray, fill=white},
  grp/.style={draw=cGray, dashed, line width=0.5pt, rounded corners=3pt, inner sep=6pt},
  grpS/.style={draw=cBlue, line width=0.8pt, rounded corners=3pt, inner sep=6pt},
  lbl/.style={font=\footnotesize, text=cGray, inner sep=1pt},
  ttl/.style={font=\small\bfseries, text=cInk, inner sep=1pt},
  ar/.style={->, draw=cInk, line width=0.6pt},
  arB/.style={->, draw=cBlue, line width=0.9pt},
  arO/.style={->, draw=cOrange, line width=0.9pt},
  arD/.style={->, draw=cGray, line width=0.6pt, dashed},
  bus/.style={draw=cGray, line width=2.4pt},
}

\begin{document}
\begin{tikzpicture}[
  u/.style={blk, font=\footnotesize, minimum height=7mm, inner sep=2pt},
  uo/.style={blkO, font=\footnotesize, minimum height=7mm, inner sep=2pt},
  ug/.style={blkG, font=\footnotesize, minimum height=7mm, inner sep=2pt},
  up/.style={blkP, font=\footnotesize, minimum height=7mm, inner sep=2pt},
]
% 左侧：平台
\node[ug, minimum width=3.2cm, minimum height=1.3cm] (cpu) at (0,0) {\textbf{E203}\\逐层命令，WFI\\PLIC 中断唤醒};
\node[ug, minimum width=3.2cm, minimum height=1.0cm] (ldr) at (0,-2.4) {\textbf{参数加载器}\\32 位 ICB 写入，门铃发布进度};
\node[up, minimum width=3.2cm] (xbar) at (0,-4.3) {系统总线仲裁器\\\scriptsize DMA > 加载器 > CPU};
\node[u, minimum width=2.0cm, minimum height=1.2cm] (b0) at (-0.6,-5.9) {SRAM 控制器\\\scriptsize 64 个 4K×16 宏};
\node[up, minimum width=1.1cm, minimum height=1.2cm] (b1) at (1.05,-5.9) {进度\\寄存器};
\node[lbl, below=1pt of b0.south east] {512 KB 共享 SRAM：单个 32 位从口，每拍 1 笔};
\draw[ar] ($(cpu.south)+(-1.2,0)$) -- node[lbl, right, pos=0.12] {CPU} ($(xbar.north)+(-1.2,0)$);
\draw[ar] (ldr.south) -- node[lbl, right] {写参数、写门铃} (xbar.north);
\draw[<->, draw=cInk, line width=0.6pt] (b0.north) -- (b0.north |- xbar.south);
\draw[<->, draw=cInk, line width=0.6pt] (b1.north) -- (b1.north |- xbar.south);
% 包装层
\node[uo, minimum width=2.4cm] (mmio) at (5.2,0.2) {MMIO 寄存器};
\node[uo, minimum width=2.4cm] (cq) at (5.2,-0.85) {命令队列（2 项）};
\node[uo, minimum width=2.4cm] (pq) at (5.2,-1.9) {参数段队列（4 项）};
\node[uo, minimum width=2.4cm, minimum height=1.0cm] (dma) at (5.2,-3.3) {流式 DMA\\\scriptsize 32 位 ICB};
\node[ttl, anchor=south] (wt) at (5.2,0.75) {\footnotesize KSE 包装层};
\begin{scope}[on background layer]
\node[draw=cOrange, line width=0.6pt, rounded corners=3pt, fit=(wt)(mmio)(dma), inner sep=5pt] (wrap) {};
\end{scope}
\draw[ar] (cpu.east) ++(0,0.2) -- node[lbl, above] {PPI MMIO} ($(mmio.west)$);
\draw[ar, draw=cOrange] (mmio.west) ++(0,-0.2) -- ++(-0.35,0) |- node[lbl, pos=0.75, below, text=cOrange!80!black] {IRQ} ($(cpu.east)+(0,-0.3)$);
\draw[ar] (mmio) -- (cq); \draw[ar] (cq) -- (pq); \draw[ar] (pq) -- (dma);
\draw[<->, draw=cInk, line width=0.6pt] (dma.south) |- node[lbl, pos=0.7, below] {读输入/参数，写输出} ($(xbar.east)$);
\draw[arD] (b1.east) -| node[lbl, pos=0.75, right] {放行上界} ($(dma.south)+(0.7,0)$);
% 引擎：左列行存储与非线性单元，右列参数通路与阵列
\node[u, minimum width=2.3cm, minimum height=1.4cm] (rsa) at (9.1,-0.6) {行存储 A\\\scriptsize 32 B × 1024 行};
\node[u, minimum width=2.3cm, minimum height=1.4cm] (rsb) at (9.1,-2.6) {行存储 B\\\scriptsize 32 B × 1024 行};
\node[uo, minimum width=1.1cm, inner sep=1pt] (smx) at (8.5,-4.2) {Softmax};
\node[uo, minimum width=1.1cm, inner sep=1pt] (ln) at (9.75,-4.2) {LN};
\node[u, minimum width=2.7cm] (pf) at (12.8,0.1) {参数预取 FIFO 2 KB};
\node[u, minimum width=2.7cm] (ld) at (12.8,-0.85) {加载器 + 驻留槽};
\node[u, minimum width=2.7cm, minimum height=1.2cm, draw=cBlue, fill=cBlueM!50] (mac) at (12.8,-2.2) {\textbf{32×32 MAC 阵列}\\\scriptsize 1024 个 s9×s8};
\node[uo, minimum width=2.7cm] (post) at (12.8,-3.6) {后处理 + GELU ROM};
\node[ug, minimum width=2.7cm] (seq) at (12.8,-4.6) {层定序器（硬连线）};
\draw[ar] (pf) -- (ld); \draw[ar] (ld) -- (mac); \draw[ar] (mac) -- (post);
\draw[ar] (rsa.east) -- (10.6,-0.6) |- node[lbl, pos=0.75, above] {读} ($(mac.west)+(0,0.3)$);
\draw[ar] ($(rsb.east)+(0,0.2)$) -- node[lbl, above, pos=0.3] {读} ($(mac.west)+(0,-0.2)$);
\draw[ar] (post.west) -- (11.0,-3.6) |- ($(rsb.east)+(0,-0.45)$);
\node[lbl, anchor=north east] at (10.95,-3.35) {写回};
\draw[<->, draw=cInk, line width=0.6pt] (smx.north) -- (smx.north |- rsb.south);
\draw[<->, draw=cInk, line width=0.6pt] (ln.north) -- (ln.north |- rsb.south);
\node[lbl, anchor=north] at (9.1,-4.6) {读写 A、B};
\draw[<->, draw=cInk, line width=0.6pt] (dma.east) ++(0,-0.25) -- (7.2,-3.55) |- (rsa.west);
\node[lbl, anchor=north west] at (6.62,-3.62) {激活};
\draw[ar] (dma.east) ++(0,0.25) -- (6.85,-3.05) -- (6.85,1.05) -| node[lbl, pos=0.25, above] {参数流} (pf.north);
\node[ttl, anchor=south east] (et) at (14.1,-5.45) {\footnotesize KSE 引擎（kse\_top）};
\begin{scope}[on background layer]
\node[draw=cBlue, line width=0.8pt, rounded corners=3pt, fit=(rsa)(pf)(seq)(smx)(et), inner sep=6pt] (eng) {};
\end{scope}
\draw[arD] (cq.east) -- (eng.west |- cq.east);
\end{tikzpicture}
\end{document}
